\subsection{The exact leading constant in odd residue characteristic}
\label{sec:odd-leading}

Write $A$ for a finite commutative local ring with residue field
$k=\mathbb F_q$ of odd characteristic, maximal ideal $J$ satisfying
$J^2=0$, cardinality $Q$, and $M=|J|=Q/q$. Multiplication makes
$J$ a vector space over $k$. Choose a representative $\widetilde\gamma$
in $A$ for each $\gamma\in k$, and put
$P_\gamma=x-\widetilde\gamma$ and
$\Lambda=\prod_{\gamma\in k}(x-\gamma)$ in $k[x]$.
Write $\overline f$ for coefficientwise reduction. Use the same bound $L\ge2$ on the sum of the two exponent gaps as in Theorem~\ref{thm:local}.
Then for each fixed odd $q$,
\[
   \lim_{Q\to\infty} qQ\bigl(1-\Delta_L(A)\bigr)=1-q^{-q}.
\]
The limit holds along any sequence of these rings with fixed residue
field size. In particular the leading constant already accounts for
the positive fraction $q^{-q}$ of residue polynomials divisible by
every monic linear polynomial.

\paragraph{Improved construction.}
At cutoff $n\ge q+3$, consider the rectangle $f_n=0$, $f_{n-1}\in J$;
subscripts indicate coefficients. Its fraction is $1/(qQ)$.
Inside this rectangle delete only the points with
$\Lambda\nmid\overline f$. Their fraction in the full ball is
$(1-q^{-q})/(qQ)$ by uniformity modulo the degree-$q$ polynomial $\Lambda$.
Also delete all points of degree at most $n-2$, and all points divisible
in $A[x]$ by some $P_\gamma^2$. The additional cost is at most
$(q+1)Q^{-2}$ since $P_\gamma$ is monic.
As in the scalar-degree gluing argument, we may remove the
zero-residue points at a cost tending to zero for this fixed ring.

Consider a remaining progression $f,y=r^i f,z=r^{i+j}f$ with a nonunit ratio.
Nonzero reductions at all three terms ensure that
$e_0=\deg\overline r$ is positive.
Write $t$ for a coefficient lift of $\overline r$ of degree $e_0$ with unit leading coefficient.
If the residue characteristic divides $j$, then $r^j=t^j$ and $y$ has
degree at most $n-je_0\le n-2$, which was excluded.
Otherwise $j$ is a unit, and the identity
$j t^{i+j} f=(i+j)t^j y-i z$ gives $\deg f\le n-ie_0$.
Thus $i=e_0=1$, and its scalar degree still obeys
$\deg\overline f\le n-(i+j)e_0\le n-2$.

Write $r=uP_\gamma+e$ with $u\in A^\times$ constant and $e\in J[x]$.
At a surviving first term $\Lambda\mid\overline f$, hence
$f=P_\gamma g+\eta$ for some $g\in A[x]$, $\eta\in J[x]$.
Thus $ef=P_\gamma eg$. For $m=i+j\ge2$,
$r^mf=u^mP_\gamma^mf+mu^{m-1}P_\gamma^{m-1}ef$ is
divisible by $P_\gamma^2$. This excludes the progression.
Passing from cutoffs to a single set by scalar-degree gluing gives
\[
 1-\Delta_L(A)\le\frac{1-q^{-q}}{qQ}+\frac{q+1}{Q^2}.
\]

\paragraph{A matching of triples.}
Start with all $f$ in the rectangle with $\Lambda\nmid a=\overline f$.
Choose the first $\beta\in k$ in a fixed order with $a(\beta)\ne0$.
Put $P=P_\beta$, $s=n-\deg a\ge2$ and
\[
 c=\frac{-f_{n-1}}{2\operatorname{lc}(a)}\in J,\qquad
 r=P+c(P^s+1),\qquad y=rf,\quad z=r^2f .
\]
The operations on $c$ use the $k$-action on $J$. These formulas give
$\deg y,\deg z\le n$ and $y_n=-c\operatorname{lc}(a)$: the coefficient
of degree $n+1$ in $z=P^2f+2cP(P^s+1)f$ is
$f_{n-1}+2c\operatorname{lc}(a)=0$.
Discard $c=0$. Each remaining first term has degree $n-1$ and $y$ has
degree $n$. Require $z_n\ne0$ as well. If $s=2$ the reduction
$\overline z=(x-\beta)^2 a$ already has degree $n$. For $s>2$,
conditional on $a$ and $c$, the coefficient $f_{n-2}$ is uniform on
$J$ and occurs in $z_n$ with coefficient one; its occurrences in terms
multiplied by $c$ depend only on its reduction. Thus these two exclusions
cost at most $2/M$ of the starts.

At $\widetilde\beta$ one has $y(\widetilde\beta)=c a(\beta)\ne0$ and
$z(\widetilde\beta)=0$. If $\gamma\ne\beta$ is a root of $a$,
then $f(\widetilde\gamma)\in J$ and
\[
 y(\widetilde\gamma)=P(\widetilde\gamma)f(\widetilde\gamma),\qquad
 z(\widetilde\gamma)=P(\widetilde\gamma)^2f(\widetilde\gamma).
\]
The factor $P(\widetilde\gamma)$ is a unit. Conditional on $a$ and $c$,
$f(\widetilde\gamma)$ is uniform on $J$: its constant coefficient has
a free uniform $J$-part. Require all these evaluations to be nonzero,
at cost at most $q/M$. At other residue points $y,z$ have nonzero
reductions. Consequently $y$ has no root among the chosen representatives,
while $z$ has exactly one, $\widetilde\beta$.

One further condition identifies the start from $y$. For every simple root
$\gamma$ of $Y=\overline y$ test the equality in $J$
\[
 y(\widetilde\gamma)=
       -y_n\,\frac{(Y/(x-\gamma))(\gamma)}{\operatorname{lc}(Y)} .
\]
It holds at $\beta$. At any other simple root of $Y$, the right side is
fixed after conditioning on $a,c$, and the left side is uniform on
$J$, as just proved. Discard starts with a second successful test,
at cost at most $q/M$. All these costs use the union bound on the
original starts and require no independence between the tests.

We obtain disjoint triples. Degree distinguishes starts from both
other positions, and evaluation at the representatives distinguishes
middles from ends. A middle identifies $\beta$ by its unique
successful test. Then $a=\overline y/(x-\beta)$ and
$c=-y_n/\operatorname{lc}(a)$ recover $r$.
An end identifies $\beta$ by its unique root and gives
$a=\overline z/(x-\beta)^2$. Dividing in $A[x]$ by the monic $P$
yields $(z/P)(\widetilde\beta)=2c a(\beta)$ and hence $c$.
Finally $r$ is a non-zero-divisor, since
$r(P-c(P^s+1))=P^2$ is monic. Each inversion thus recovers at most one
start. The ordinary triple belongs to the allowed family for each $L\ge2$. Removing one member of each disjoint triple forces
\[
 1-\Delta_L(A)\ge
 \frac{1-q^{-q}}{qQ}
 \left(1-\frac{2+2q}{M}\right).
\]
Together the two bounds give the limit.
